\textbf{\textit{Soal.}} Diberikan $\triangle ABC$ dimana $A',B',C'$ berturut-turut adalah pencerminan $A,B,C$ terhadap $BC,CA,AB$. Perpotongan lingkaran luar $\triangle ABB'$ dan $\triangle ACC'$ adalah $A_1$. Definisikan $B_1$ dan $C_1$ secara serupa. Buktikan bahwa $AA_1,BB_1,$ dan $CC_1$ konkuren (bertemu di satu titik).

\begin{proof}[\textbf{Solusi. }] (Wednesday, 6 August 2025)     
    Misalkan $AC$ berpotongan dengan lingkaran $(ABB')$ untuk kedua kalinya di $A_B$ dan misalkan $AB$ berpotongan dengan lingkaran $(ACC')$ untuk kedua kalinya di $A_C$. Perhatikan karena $B'$ adalah hasil pencerminan $B$ terhadap $AA_B$, maka $AA_B$ adalah diameter lingkaran $(ABA_BB')$ sehingga  $BA_B \perp AB$. Dengan cara yang serupa dapat diperoleh bahwa juga bahwa $AA_C$ adalah diameter lingkaran $(ACA_CC')$ sehingga $CA_C \perp AC$.
    
    Dari fakta-fakta tersebut, didapat $\angle A_BA_1A = \angle AA_1A_C = 90^\circ$, maka $A_B, A_1, A_C$ segaris dan juga didapat bahwa $AA_1$ adalah garis tinggi $\triangle AA_BA_C$.
    
    Sekarang, misalkan $H_A$ adalah perpotongan antara $CA_C$ dan $BA_B$, maka $\angle ACH_A = \angle H_ABA = 90^\circ$ yang mengakibatkan $H_A$ adalah titik tinggi $\triangle AA_CA_B$ sekaligus $AH_A$ menjadi diameter lingkaran luar $\triangle ABC$. Karena $AA_1$ adalah garis tinggi $\triangle AA_BA_C$, maka $H_A$ berada di garis $AA_1$.
    
    Definisikan $H_B$ dan $H_C$ secara serupa dengan definisi $H_A$, maka didapat pula bahwa $H_B$ dan $H_C$ berturut-turut berada di garis $BB_1$ dan $CC_1$. Oleh karena itu, $AH_A$, $BH_B$, dan $CH_C$ adalah diameter lingkaran luar $\triangle ABC$, maka didapat $AH_A$, $BH_B$, dan $CH_C$ berpotongan di satu titik yaitu pusat lingkaran. Karena $A_1,B_1,C_1$ berturut-turut ada di garis $AH_A$, $BH_B$, $CH_C$ maka hal ini berakibat $AA_1$, $BB_1$, dan $CC_1$ berpotongan di satu titik. \qed
\end{proof}